% Part: proof-theory
% Chapter: cut-elimination
% Section: ce-largest

\documentclass[../../../include/open-logic-section]{subfiles}

\begin{document}

\olfileid{pt}{cut}{inv}

\olsection{移除最大切割}

如 \olref[seq][inv]{lem:G3c-invert} 的证明所示，若 \Log{G3c}
推导出某逻辑规则（除
\RightR{\lexists} 和 \LeftR{\lforall} 外）的结论，则它也能推导出该规则的每个前提，且证明高度不会增加。
我们可以进一步证明，该性质对于 $\Log{G3c} + \CutCS$ 也成立。

和之前一样，我们将一个 \CutCS{} 推理的\emph{切割秩}定义为其切割公式的深度。

\begin{lem}\ollabel{lem:inv-G3c-cut}
如果 $\Log{G3c} + \CutCS$ 使用推导~$\pi$ 推出了某个逻辑规则~$R$（\RightR{\lexists} 或 \LeftR{\lforall} 除外）的结论，那么它也能使用推导~$\pi'$（或根据该规则前提的数量，使用推导~$\pi'$、$\pi''$）推出该规则的每个前提。此外，
(a)~$\pheight{\pi'}\le\pheight{\pi}$（及
$\pheight{\pi''}\le\pheight{\pi}$），且 (b)~$\pi'$（以及 $\pi''$）中 \CutCS{}
推理的个数及其秩与~$\pi$ 中的相同。
\end{lem}

\begin{proof}
通过检视
\cref{pt:seq:inv:lem:G3c-invert,pt:seq:inv:lem:invert-quant} 的证明即可得到，该证明是对~$\pheight{\pi}$ 进行归纳。在基始情况中，我们将一个公理替换为另一个公理。因此条件 (a) 和~(b) 平凡成立。如果~$\pi$ 的最后一条规则就是~$R$，那么它的直接子推导即为所需的推导~$\pi_i$，并且 (a) 和 (b) 都成立。在所有其他情况下，我们对~$\pi$ 的直接子推导（即导出最后一条推理~$R$ 之前提的那些推导）应用归纳假设，然后将同样的规则 $R$ 应用于所得结果。条件 (a) 和 (b) 对于新证明的直接子推导成立，因此 (a) 和 (b) 对整个证明也成立。余下只需验证最后一条推理是~\CutCS{} 的情况。我们再次仅通过一个例子来说明：假设 $\pi$ 是 \LeftR{\land} 的结论 $!B \land !C, \Gamma \Sequent \Delta$ 的一个推导，并且以~\CutCS{} 结束：
\[
\AxiomC{}
\RightLabel{$\pi_1$}
\Deduce$!B \land !C, \Gamma \fCenter \Delta, !A$
\AxiomC{}
\RightLabel{$\pi_2$}
\Deduce$!A, !B \land !C, \Gamma \fCenter \Delta$
\RightLabel{\CutCS}
\BinaryInf$!B \land !C, \Gamma \fCenter \Delta$
\DisplayProof
\]
对 $\pi_1$ 和~$\pi_2$（它们也是 $\LeftR{\land}$ 结论之实例的推导）应用归纳假设，可分别得到 $!B, !C, \Gamma \fCenter
\Delta, !A$ 和 $!A, !B, !C, \Gamma \fCenter \Delta$ 的推导 $\pi_1'$ 和~$\pi_2'$。我们利用 \CutCS{} 将它们组合起来得到~$\pi'$：
\[
\AxiomC{}
\RightLabel{$\pi_1'$}
\Deduce$!B, !C, \Gamma \fCenter \Delta, !A$
\AxiomC{}
\RightLabel{$\pi_2'$}
\Deduce$!A, !B, !C, \Gamma \fCenter \Delta$
\RightLabel{\CutCS}
\BinaryInf$!B, !C, \Gamma \fCenter \Delta$
\DisplayProof
\]
显然，(a) 和 (b) 是满足的。
\end{proof}

请注意，如果我们使用 \Cut{} 而不是 \CutCS，这将不可行。因为那样的话，最后那个推导的结论中，前件会有两份 $!B, !C, \Gamma$ 而后件会有两份 $\Delta$。我们就需要诉诸收缩规则的可容许性，但 \olref[seq][inv]{prop:G3c-cont-adm} 的证明又用到了可逆性引理。这样一来我们就会陷入循环论证。

我们可以利用 \cref{pt:cut:inv:lem:inv-G3c-cut} 来给出切割消去的一个替代证明。在这个证明中，我们不是逐次移除最顶层的 \CutCS{} 推理，而是移除具有最大秩的 \CutCS{} 推理。也就是说，我们从 $\Log{G3c} + \CutCS$ 中的一个推导~$\pi$ 开始，移除其中的某个切割（可能用秩更小的切割替换它）——然后如此反复，直到不含任何切割为止。唯一的条件是，在我们所处理的切割之上，不存在秩相同或更高的其他切割。

\begin{lem}\ollabel{lem:max-cut-red-G3c} 如果 $\pi$ 是 $\Log{G3c}+\CutCS$ 中的一个推导，它以秩为~$n$ 的一个~\CutCS{} 推理结束，并且除此之外仅使用了~$\Log{G3c}$ 的规则以及秩小于~$n$ 的 \CutCS{} 推理，那么在~$\Log{G3c} + \CutCS$ 中存在同一末相继式的证明~$\pi'$，其中每个 \CutCS{} 推理的秩都小于~$n$。
\end{lem}

\begin{proof}
推导~$\pi$ 结束于：
\[
\AxiomC{}
\RightLabel{$\pi_1$}
\Deduce$\Gamma \fCenter \Delta, !A$
\AxiomC{}
\RightLabel{$\pi_2$}
\Deduce$!A, \Gamma \fCenter \Delta$
\RightLabel{\CutCS}
\BinaryInf$\Gamma \fCenter \Delta$
\DisplayProof
\]
我们根据~$!A$ 的形式区分情况。

假设 $!A$ 是原子的。注意到，由于~$\pi_1$ 和 $\pi_2$ 中所有 \CutCS{} 推理的秩必须 $< \depth{!A} = 0$，可知~$\pi_1$ 和~$\pi_2$ 中不可能存在 \CutCS{} 推理。我们通过对 $\pheight{\pi_2}$ 进行归纳来证明 $\Gamma \Sequent \Delta$ 有一个无切割的证明。若 $\pheight{\pi_2} = 0$，则 $!A, \Gamma \Sequent \Delta$ 是一条公理。如果 $!A$ 不是主公式，则某个原子公式 $!C$ 同时属于~$\Gamma$ 和~$\Delta$，于是 $\Gamma \Sequent \Delta$ 本身便是公理。若 $!A$ 是主公式，则 $\Delta = \Delta, !A$。此时，$\pi_1$ 结束于 $\Gamma \Sequent \Delta', !A, !A$。我们应用 \olref[seq][inv]{prop:G3c-cont-adm}（收缩规则的可容许性）即可得到 $\Gamma \Sequent \Delta', !A$ 的一个无切割证明。若 $\pheight{\pi_2} > 0$，则 $\pi_2$ 结束于某个规则~$R$。取该规则的一个前提：它具有形式 $!A, \Gamma',
\Pi \Sequent \Delta', \Lambda$，其中 $\Pi$ 和 $\Lambda$ 是~$R$ 的活动公式。将 \olref[seq][adm]{prop:weak-G3c-adm} 应用于~$\pi_1$ 和 $\pi_2$，可得到 $\Gamma, \Gamma',
\Pi \Sequent \Delta, \Delta', \Lambda, !A$ 和 $!A, \Gamma, \Gamma',
\Pi \Sequent \Delta, \Delta', \Lambda$ 的推导，进而可应用归纳假设。由此便给出了 $R$ 的每个前提 $\Gamma, \Gamma', \Pi
\Sequent \Delta, \Delta', \Lambda$ 的推导。应用 $R$ 得到 $\Gamma, \Gamma \Sequent \Delta, \Delta$，再运用 \olref[seq][inv]{prop:G3c-cont-adm} 即得到 $\Gamma \Sequent \Delta$ 的无切割证明。

如果 $!A$ 不是原子的，我们根据~$!A$ 的主联结词区分情况。在每一种情况下，我们应用 \olref{lem:inv-G3c-cut} 来获得以~$!A$ 为主公式的左规则和右规则前提的推导。然后我们按照 \olref[top]{lem:cut-adm-G3c} 证明中情况~E 的方式处理。
\begin{enumerate}
  \item $!A \ident \lnot !B$。推导 $\pi_1$ 和 $\pi_2$ 结束于 $\Gamma \Sequent \Delta, \lnot !B$ 和 $\lnot !B, \Gamma \Sequent \Delta$。由 \olref{lem:inv-G3c-cut}，存在 $!B, \Gamma \Sequent \Delta$ 和 $\Gamma \Sequent \Delta, !B$ 的推导 $\pi_1'$ 和 $\pi_2'$。推导~$\pi'$ 为：
  \[
  \AxiomC{}
  \RightLabel{$\pi_2'$}
  \Deduce$\Gamma \fCenter \Delta, !B$
  \AxiomC{}
  \RightLabel{$\pi_1'$}
  \Deduce$!B, \Gamma \fCenter \Delta$
  \RightLabel{\CutCS}
  \BinaryInf$\Gamma \fCenter \Delta$
  \DisplayProof
  \]
  \item $!A \ident (!B \lor !C)$。推导 $\pi_1$ 和 $\pi_2$ 结束于 $\Gamma \Sequent \Delta, !B \lor !C$ 和 $!B \lor !C, \Gamma \Sequent \Delta$。由 \olref{lem:inv-G3c-cut}（对~$\pi_1$ 应用 \RightR{\lor} 的可逆性），存在 $\Gamma \Sequent \Delta, !B, !C$ 的推导 $\pi_1'$。由 \olref{lem:inv-G3c-cut}（对~$\pi_2$ 应用 \LeftR{\lor} 的可逆性），存在 $!B, \Gamma \Sequent \Delta$ 的推导 $\pi_2'$ 和 $!C, \Gamma \Sequent \Delta$ 的推导 $\pi_2''$。将 \olref[seq][adm]{prop:weak-G3c-adm} 应用于 $\pi_2''$，我们还得到 $!C, \Gamma \Sequent \Delta, !B$ 的推导~$\pi_2'''$。推导~$\pi'$ 为：
  \[
  \AxiomC{}
  \RightLabel{$\pi_1'$}
  \Deduce$\Gamma \fCenter \Delta, !B, !C$
  \AxiomC{}
  \RightLabel{$\pi_2'''$}
  \Deduce$!C, \Gamma \fCenter \Delta, !B$
  \RightLabel{\CutCS}
  \BinaryInf$\Gamma \fCenter \Delta, !B$
  \AxiomC{}
  \RightLabel{$\pi_2'$}
  \Deduce$!B, \Gamma \fCenter \Delta$
  \RightLabel{\CutCS}
  \BinaryInf$\Gamma \fCenter \Delta$
  \DisplayProof
  \]
  \item $!A \ident (!B \land !C)$。练习。
  \item $!A \ident (!B \lif !C)$。推导 $\pi_1$ 和 $\pi_2$ 结束于 $\Gamma \Sequent \Delta, !B \lif !C$ 和 $!B \lif !C, \Gamma \Sequent \Delta$。由 \olref{lem:inv-G3c-cut}（对~$\pi_1$ 应用 \RightR{\lif} 的可逆性），存在 $!B, \Gamma \Sequent \Delta, !C$ 的推导 $\pi_1'$。由 \olref{lem:inv-G3c-cut}（对~$\pi_2$ 应用 \LeftR{\lif} 的可逆性），存在 $\Gamma \Sequent \Delta, !B$ 的推导 $\pi_2'$ 和 $!C, \Gamma \Sequent \Delta$ 的推导 $\pi_2''$。将 \olref[seq][adm]{prop:weak-G3c-adm} 应用于 $\pi_2''$，我们还得到 $!C, !B, \Gamma \Sequent \Delta$ 的推导~$\pi_2'''$。推导~$\pi'$ 为：
  \[
  \AxiomC{}
  \RightLabel{$\pi_2'$}
  \Deduce$\Gamma \fCenter \Delta, !B$
  \AxiomC{}
  \RightLabel{$\pi_1'$}
  \Deduce$!B, \Gamma \fCenter \Delta, !C$
  \AxiomC{}
  \RightLabel{$\pi_2'''$}
  \Deduce$!C, !B, \Gamma \fCenter \Delta$
  \RightLabel{\CutCS}
  \BinaryInf$!B, \Gamma \fCenter \Delta$
  \RightLabel{\CutCS}
  \BinaryInf$\Gamma \fCenter \Delta$
  \DisplayProof
  \]
  \item $!A \ident \lforall[x][!B(x)]$。推导 $\pi_1$ 和 $\pi_2$ 结束于 $\Gamma \Sequent \Delta, \lforall[x][!B(x)]$ 和 $\lforall[x][!B(x)], \Gamma \Sequent \Delta$。由 \olref{lem:inv-G3c-cut}，存在 $\Gamma \Sequent \Delta, !B(c)$ 的推导~$\pi_1'$。

  现在考虑推导~$\pi_2$。它结束于 $\lforall[x][!B(x)],
  \Gamma \Sequent \Delta$。从~$\pi_2$ 的末相继式向上追溯，标记相继式左侧每一个“导向”~$\pi_2$ 末相继式中 $\lforall[x][!B(x)]$ 之出现的 $\lforall[x][!B(x)]$ 的出现，即：
  (a)~在~$\pi_2$ 的末相继式中标记 $\lforall[x][!B(x)]$。(b)~如果 $\lforall[x][!B(x)]$ 出现在某条规则结论的上下文中且已被标记，则在该规则的前提中标记对应的出现。(c)~如果 $\lforall[x][!B(x)]$ 是 \LeftR{\lforall} 规则结论中的主公式且已被标记，则在该规则的前提中标记对应的 $\lforall[x][!B(x)]$ 及 $!B(t)$ 的活动出现。

  现在考虑所有这些被标记的 $\lforall[x][!B(x)]$ 出现。
  它们均未在 \LeftR{\lforall} 以外的规则中作为活动公式出现（只有 \LeftR{\lforall} 的活动公式会被标记）。在 $\pi_2$ 中删除所有被标记的 $\lforall[x][!B(x)]$ 出现。
  若删除后仍为正确的推导，则说明被标记的 $\lforall[x][!B(x)]$ 从未作为活动公式出现；它们都直接源自公理。该推导结束于~$\Gamma \Sequent \Delta$，我们可以用它替换 $\pi$。由此我们便移除了一个最大秩的切割。

  否则，所得便不是正确的证明，因为我们删除了在某些 \LeftR{\lforall} 推理中作为活动公式的 $\lforall[x][!B(x)]$。我们将这些推理转变为了如下形式的“推理”：
  \[
  \AxiomC{}
  \RightLabel{$\pi_t$}
  \Deduce$!B(t), \Pi \fCenter \Lambda$
  \RightLabel{\LeftR{\lforall}}
  \UnaryInf$\Pi \fCenter \Lambda$
  \DisplayProof
  \]
  将每一个位于最顶层的此类“推理”替换为
  \[
  \AxiomC{}
  \RightLabel{$\Subst{\pi_1'}{t}{c}[\Pi \Sequent \Lambda]$}
  \Deduce$\Gamma, \Pi \fCenter \Delta, \Lambda, !B(t)$
  \AxiomC{}
  \RightLabel{$\pi_t[\Gamma \Sequent \Delta]$}
  \Deduce$!B(t), \Gamma, \Pi \fCenter \Delta, \Lambda$
  \RightLabel{\CutCS}
  \BinaryInf$\Gamma, \Pi \fCenter \Delta, \Lambda$
  \DisplayProof
  \]
  并在其下的每个相继式左边添加 $\Gamma$，右边添加 $\Delta$。重复此过程，直到所有前提中含有被标记的 $!B(t)$ 的 \LeftR{\lforall} “推理”都被替换为关于某个~$!B(t)$ 的 \CutCS{} 推理。所得推导（此时它已是正确的推导）的末相继式形如 $\Gamma, \dots, \Gamma \Sequent \Delta, \dots, \Delta$。应用收缩规则的可容许性将其变为~$\Gamma \Sequent \Delta$ 的推导，并用它替换 $\pi$。这样我们就把一个关于 $\lforall[x][!B(x)]$ 的切割替换为（可能多个）关于 $!B(t)$ 的切割，而这些切割的秩均更低。任何原先在 $\pi_1$ 或 $\pi_2$ 中已存在的切割仍然保留，但它们的秩也都更低。
\end{enumerate}
\end{proof}

\begin{prob}
请在 \olref{lem:inv-G3c-cut} 的证明中验证 $!A \ident (!B \land !C)$ 的情形。也就是说，证明：如果存在推导~$\pi$ 以
\[
\AxiomC{}
\RightLabel{$\pi_1$}
\Deduce$\Gamma \fCenter \Delta, !B \land !C$
\AxiomC{}
\RightLabel{$\pi_2$}
\Deduce$!B \land !C, \Gamma \fCenter \Delta$
\RightLabel{\CutCS}
\BinaryInf$\Gamma \fCenter \Delta$
\DisplayProof
\]
结尾，且 $\pi_1$ 和 $\pi_2$ 仅包含秩 $<\depth{!B \land !C}$ 的切割，那么存在 $\Gamma \Sequent \Delta$ 的推导~$\pi'$，且 $\pi'$ 仅包含秩 $<\depth{!B \land !C}$ 的切割。
\end{prob}

% \begin{prob}
% Verify the case where $!A \ident \lexists[x][!B(x)]$ in the proof of
% \olref{lem:inv-G3c-cut}.
% \end{prob}

\olref[top]{lem:cut-adm-G3c} 表明，我们可以将推导中秩最高的 \CutCS{} 推理替换为秩更低的 \CutCS{} 推理。我们可以再次利用该引理，通过逐次将其应用于以最高秩 \CutCS{} 结尾的最顶层子推导，来从一个推导中消去任意数量的 \CutCS{} 推理。

\begin{thm}\ollabel{thm:cut-elim-G3c-largest}
$\Log{G3c} + \CutCS$ 中的任意推导 $\pi$ 都可转化为 \Log{G3c} 中的证明。
\end{thm}

\begin{proof}
  首先证明，可以消去 $\pi$ 中所有秩最大的切割。设 $d$ 为 $\pi$ 中出现的切割的最大秩，并设 $n$ 为秩为 $d$ 的切割的数量。下对 $n$ 进行归纳。若 $n=0$，则无需证明。若 $n>0$，选取一个秩为 $d$ 的最顶层切割，并令 $\pi_1$ 为以此切割结尾的子推导。根据 \olref{lem:max-cut-red-G3c}，存在一个具有相同末端相继式的推导~$\pi_1'$，其中所有切割的秩均 $<d$。将 $\pi$ 中的 $\pi_1$ 替换为 $\pi_1'$。由此得到一个包含 $n-1$ 个秩为 $d$ 的切割的推导，从而可以应用归纳假设。

  接下来，对 $d$ 进行归纳即可得证。如果 $d=0$，所有切割都是原子的，由前述结果可知它们均可被消去。如果 $d>0$，前述结果可得到一个证明，其中所有秩为 $d$ 的切割已被秩 $<d$ 的切割替换。由于 $d$ 是 $\pi$ 中切割的最大秩，我们便得到了一个切割最大秩 $<d$ 的证明，从而可以应用归纳假设。
\end{proof}

\end{document}